\documentclass[12pt]{article}

% Language setting
\usepackage[brazilian]{babel}
\usepackage[utf8]{inputenc}
\usepackage[T1]{fontenc}

% Fonte Times New Roman
\usepackage{mathptmx}

% Espaçamento 1.5 linhas
\usepackage{setspace}
\onehalfspacing

% Margens: 3cm esquerda/superior, 2cm direita/inferior
\usepackage[a4paper,top=3cm,bottom=2cm,left=3cm,right=2cm]{geometry}

% Useful packages
\usepackage{amsmath}
\usepackage{amssymb}
\usepackage{graphicx}
\usepackage[colorlinks=true, allcolors=blue]{hyperref}
\usepackage{booktabs}
\usepackage{longtable}
\usepackage{array}
\usepackage{multirow}
\usepackage{float}
\usepackage{enumitem}
\usepackage{xcolor}
\usepackage{titlesec}

% Apenas números de página no rodapé
\pagestyle{plain}

% Title formatting
\titleformat{\section}{\Large\bfseries}{\thesection}{1em}{}
\titleformat{\subsection}{\large\bfseries}{\thesubsection}{1em}{}
\titleformat{\subsubsection}{\normalsize\bfseries}{\thesubsubsection}{1em}{}

% Comando para links textuais às referências bibliográficas
\DeclareRobustCommand{\citelink}[2]{\hyperlink{cite.#1}{#2}}

\begin{document}

\begin{center}
    Universidade de São Paulo \\
    Instituto de Ciências Matemáticas e de Computação \\
    \vspace{0.8cm}
    {\Large \textbf{Métodos de Solução para o Agendamento e Roteamento de Equipes para a Assistência Domiciliar}} \\
\end{center}

\vspace{0.3cm}
\noindent
Aluno: Fernando Alee Suaiden \\
Orientadora: Maristela Oliveira dos Santos \\
Relatório Semestral -- 2 de março de 2026 \\
Processo FAPESP: 2025/21835-0

\vspace{0.5cm}
\noindent
\textbf{RESUMO:} A otimização de rotas e agendamentos das equipes de Atenção Domiciliar do Programa Melhor em Casa é modelada como um Home Health Care Routing and Scheduling Problem (HHC-RSP). Foi realizada uma Análise Exploratória de Dados utilizando bases públicas do CNES/DATASUS, CBO e IBGE, permitindo caracterizar o cenário operacional do programa e extrair os parâmetros fundamentais para modelagem. Foram identificadas 2.664 equipes ativas em apenas 21,9\% dos municípios brasileiros, com distribuição desigual entre regiões e composição heterogênea por tipo de equipe. A análise de conformidade legal revelou que 76,7\% das equipes atendem à Portaria GM/MS nº 3.005/2024, mas a EMAD~I, tipo mais numeroso, apresenta apenas 64,6\% de conformidade, com a maioria dos déficits concentrados no requisito de enfermagem. Foram quantificados os parâmetros de capacidade ($Q_k$) e habilidades ($S_k$) de cada equipe e geradas instâncias sintéticas para testes do algoritmo BRKGA.

\vspace{0.3cm}
\noindent
\textbf{Palavras-chave:} Atenção Domiciliar; Roteamento de Veículos; HHC-RSP; BRKGA; Programa Melhor em Casa; Otimização Combinatória; SUS.

\newpage

\section{Contextualização e Proposta}

\subsection{Contexto e Motivação}

O Brasil vive uma acelerada transição demográfica. Segundo o \citelink{ibge2022}{Censo 2022 do IBGE}, a população com 60 anos ou mais atingiu 32,1 milhões de pessoas (15,6\% do total), um crescimento de 56\% em relação a 2010, quando representava 10,8\%. O índice de envelhecimento saltou de 44,8 para 80,0 no período, indicando que há atualmente 80 idosos para cada 100 crianças de 0 a 14 anos. Em estados como Rio Grande do Sul e Rio de Janeiro, o número de idosos já supera o de crianças. Este fenômeno demográfico pressiona diretamente o sistema de saúde, aumentando a demanda por cuidados de longa duração e atendimento domiciliar.

Nesse contexto, o \citelink{portaria825}{Programa Melhor em Casa}, instituído pelo Ministério da Saúde em 2011, oferece atendimento domiciliar por meio de equipes multiprofissionais que realizam visitas regulares a pacientes que necessitam de cuidados contínuos, mas que podem ser tratados em casa.

\subsection{O Problema de Otimização (HHC-RSP)}

O Home Health Care Routing and Scheduling Problem (HHC-RSP) é um problema de otimização combinatória que visa determinar as melhores rotas e agendamentos para as equipes de atendimento domiciliar. Pertence à família dos Vehicle Routing Problems (VRP) com características específicas do setor de saúde. Dados um conjunto de pacientes com localização, janelas de tempo e requisitos de habilidades ($R_i$), e um conjunto de equipes com base de operação, capacidade ($Q_k$) e habilidades ($S_k$), o problema busca encontrar rotas e agendamentos que minimizem custos de deslocamento e tempo de espera, sujeito a restrições de capacidade, habilidades, janelas de tempo e cobertura.

\subsection{Objetivo e Metodologia}

O objetivo deste projeto é estudar métodos de otimização para a determinação de rotas e agendamentos das equipes de Atenção Domiciliar no contexto do \citelink{portaria825}{Programa Melhor em Casa}. Para isso, utiliza dados públicos do \citelink{cnes}{CNES/DATASUS} (Cadastro Nacional de Estabelecimentos de Saúde), CBO (Classificação Brasileira de Ocupações) e \citelink{ibge2022}{IBGE} para caracterizar as equipes e gerar instâncias realísticas do HHC-RSP. O método de solução adotado é o BRKGA (Biased Random-Key Genetic Algorithm), seguindo a formulação de \citelink{mankowska2014}{Mankowska et al. (2014)} e a abordagem computacional de \citelink{kummer2024}{Kummer et al. (2024)}.

\section{Realizações no Período}

Este relatório cobre o primeiro período de desenvolvimento (outubro/2025 a março/2026), correspondente à revisão bibliográfica e à compreensão do problema no contexto brasileiro.

\subsection{Revisão Bibliográfica}

Foi realizada revisão da literatura sobre VRP, HHC-RSP e métodos de solução. O trabalho seminal de \citelink{dantzig1959}{Dantzig e Ramser (1959)} introduziu o Vehicle Routing Problem, estabelecendo as bases para toda a família de problemas de roteamento. \citelink{mankowska2014}{Mankowska et al. (2014)} propuseram o modelo de referência para HHC-RSP com serviços interdependentes, considerando sincronização entre profissionais, e \citelink{fikar2017}{Fikar e Hirsch (2017)} realizaram uma revisão sistemática da literatura, classificando 41 artigos por características do problema e métodos de solução. \citelink{goncalves2011}{Gonçalves e Resende (2011)} introduziram o BRKGA, meta-heurística evolucionária baseada em chaves aleatórias, posteriormente aplicada ao HHC-RSP brasileiro por \citelink{kummer2024}{Kummer et al. (2024)}, com resultados computacionais em instâncias de médio porte.

\subsection{Análise Exploratória de Dados -- Programa Melhor em Casa}

O principal resultado deste período foi uma Análise Exploratória de Dados (AED) completa das bases públicas do governo. A modelagem do HHC-RSP requer parâmetros realísticos, como capacidade ($Q_k$), habilidades ($S_k$) e localização das equipes, que só podem ser obtidos a partir de uma caracterização detalhada do programa. A análise segue um percurso do panorama nacional à escala estadual, examinando sucessivamente a distribuição geográfica, a composição por tipo de equipe e a conformidade com a legislação vigente.

\subsubsection{Fontes de Dados e Metodologia}

Foram utilizadas três fontes de dados públicas: o \citelink{cnes}{CNES/DATASUS} (competência 08/2025), com informações sobre estabelecimentos de saúde, equipes e profissionais; a Classificação Brasileira de Ocupações (CBO, edição 2024); e dados demográficos e geográficos do \citelink{ibge2022}{Censo 2022 do IBGE}. A partir dessas bases, as equipes de Atenção Domiciliar foram identificadas e filtradas, seus profissionais foram vinculados via tabelas relacionais, e as métricas de capacidade ($Q_k$), habilidades ($S_k$) e conformidade legal foram calculadas para cada equipe. Os detalhes técnicos do processamento estão documentados no repositório do projeto.

\subsubsection{Panorama Nacional}

O Brasil possui 2.664 equipes ativas de Atenção Domiciliar, organizadas em quatro tipos:

\begin{table}[H]
\centering
\begin{tabular}{|l|c|c|l|}
\hline
\textbf{Tipo de Equipe} & \textbf{Quantidade} & \textbf{Percentual} & \textbf{Função} \\
\hline
EMAD I (código 22) & 1.093 & 41,0\% & Atendimento direto \\
EMAP (código 23) & 929 & 34,9\% & Apoio especializado \\
EMAD II (código 46) & 460 & 17,3\% & Atendimento (menor porte) \\
EMAP-R (código 77) & 182 & 6,8\% & Apoio rural/reabilitação \\
\hline
\textbf{TOTAL} & \textbf{2.664} & \textbf{100\%} & -- \\
\hline
\end{tabular}
\caption{Distribuição nacional das equipes por tipo.}
\end{table}

As EMAD (Equipes Multiprofissionais de Atenção Domiciliar, tipos I e II) são responsáveis pelo atendimento direto nas modalidades AD2 e AD3 (níveis intermediário e avançado de cuidado domiciliar), definidas pela \citelink{portaria825}{Portaria GM/MS nº 825/2016}. As EMAP (Equipes Multiprofissionais de Apoio) fornecem suporte especializado. A diferença entre os dois tipos de EMAD está nos requisitos de composição: a EMAD~I exige maior Carga Horária Semanal (CHS) de médicos e enfermeiros, sendo voltada a municípios de maior porte, enquanto a EMAD~II é uma versão mais leve, adequada a municípios menores. Da mesma forma, a EMAP-R possui requisitos reduzidos em relação à EMAP convencional, sendo destinada a municípios com menos de 20 mil habitantes.

Essas 2.664 equipes estão presentes em apenas 1.218 dos 5.570 municípios brasileiros, correspondendo a uma cobertura de 21,9\%. Essa cobertura limitada tem consequência direta para o HHC-RSP: como os recursos são escassos, a otimização das rotas e agendamentos torna-se necessária para maximizar o número de pacientes atendidos com a capacidade disponível. Além disso, os quatro tipos de equipe possuem capacidades ($Q_k$) e habilidades ($S_k$) distintas, o que configura um problema de frota heterogênea.

\subsubsection{Disparidades Regionais}

A média nacional de 21,9\% oculta diferenças substanciais entre as regiões do país.

\begin{table}[H]
\centering
\begin{tabular}{|l|c|c|c|c|}
\hline
\textbf{Região} & \textbf{Mun. Cobertos} & \textbf{Cobertura (\%)} & \textbf{Equipes} & \textbf{Eq./100k hab.} \\
\hline
Nordeste & 564 & 31,4\% & 1.104 & 2,02 \\
Norte & 99 & 22,0\% & 237 & 1,35 \\
Sudeste & 359 & 21,5\% & 892 & 1,05 \\
Centro-Oeste & 89 & 19,1\% & 216 & 1,33 \\
Sul & 107 & 9,0\% & 215 & 0,72 \\
\hline
\textbf{Brasil} & \textbf{1.218} & \textbf{21,9\%} & \textbf{2.664} & \textbf{1,31} \\
\hline
\end{tabular}
\caption{Cobertura do Programa Melhor em Casa por região.}
\end{table}

A tabela revela diferenças expressivas entre regiões. O Nordeste lidera em cobertura, enquanto o Sul apresenta a menor penetração do programa. Para o HHC-RSP, essa disparidade implica que a estrutura do problema varia por localidade: regiões com poucas equipes e grande dispersão geográfica geram instâncias com deslocamentos mais longos e menor flexibilidade na alocação.

Além da distribuição geográfica, as regiões diferem na composição das equipes implantadas.

\begin{table}[H]
\centering
\begin{tabular}{|l|c|c|c|c|}
\hline
\textbf{Região} & \textbf{EMAD I} & \textbf{EMAD II} & \textbf{EMAP} & \textbf{EMAP-R} \\
\hline
Nordeste (1.104 eq.) & 27,9\% & 23,8\% & 37,8\% & 10,5\% \\
Norte (237 eq.) & 45,6\% & 13,1\% & 40,1\% & 1,3\% \\
Centro-Oeste (216 eq.) & 38,4\% & 19,9\% & 40,3\% & 1,4\% \\
Sudeste (892 eq.) & 53,0\% & 9,8\% & 31,2\% & 6,1\% \\
Sul (215 eq.) & 56,3\% & 16,7\% & 24,2\% & 2,8\% \\
\hline
\textbf{Brasil} (2.664 eq.) & 41,0\% & 17,3\% & 34,9\% & 6,8\% \\
\hline
\end{tabular}
\caption{Composição percentual por tipo de equipe em cada região.}
\end{table}

O Nordeste apresenta um perfil de composição distinto das demais regiões, com maior presença de EMAD~II e EMAP-R, tipos com requisitos menos exigentes. As EMAD~II possuem metade da carga horária médica e de enfermagem em relação às EMAD~I, enquanto as EMAP-R exigem CHS total de apenas 60 horas (contra 90 horas das EMAP convencionais). O Sudeste e o Sul, por outro lado, concentram mais da metade de suas equipes no tipo EMAD~I.

Essa diferença de composição é relevante tanto para a conformidade legal, pois cada tipo de equipe possui requisitos distintos de CHS e categorias profissionais, quanto para a otimização, pois define a heterogeneidade da frota em termos de capacidade ($Q_k$) e habilidades ($S_k$).

\subsubsection{Conformidade com a Portaria GM/MS nº 3.005/2024}

A caracterização por tipo de equipe torna possível a verificação de conformidade com os requisitos legais de composição. A \citelink{portaria3005}{Portaria GM/MS nº 3.005/2024} estabelece requisitos mínimos para cada tipo de equipe. As EMAD~I exigem no mínimo 40h de médico, 60h de enfermeiro, 120h de técnico de enfermagem e 30h de fisioterapeuta ou assistente social. As EMAD~II possuem requisitos reduzidos pela metade para médico (20h) e enfermeiro (30h). As EMAP e EMAP-R devem ter pelo menos três categorias de nível superior, com CHS total mínima de 90h e 60h, respectivamente. Além disso, nenhum profissional de EMAD pode ter CHS inferior a 20h (Art. 547, §1º).

A verificação consistiu em agregar as CHS por categoria profissional (via prefixos CBO) para cada equipe, descartar profissionais com CHS inferior a 20h, e comparar os valores resultantes com os requisitos mínimos. Os resultados, ilustrados na Figura~\ref{fig:conformidade_tipo}, revelam que 76,7\% das equipes estão em conformidade, mas essa média esconde uma diferença importante.

\begin{figure}[H]
\centering
\includegraphics[width=0.85\textwidth]{conformidade_por_tipo.png}
\caption{Conformidade legal por tipo de equipe (\citelink{portaria3005}{Portaria GM/MS nº 3.005/2024}).}
\label{fig:conformidade_tipo}
\end{figure}

A EMAD~I, tipo mais numeroso e central para o atendimento direto, é justamente o que apresenta a menor taxa de conformidade (64,6\%). As EMAD~II e EMAP, com requisitos menos exigentes, atingem taxas superiores a 86\%.

A decomposição das 387 EMAD~I não conformes revela uma causa predominante: 298 equipes (77,0\% das não conformes) possuem carga horária de enfermeiro inferior a 60h, seguidas por déficit de técnico de enfermagem em 84 equipes (21,7\%), de fisioterapeuta ou assistente social em 33 (8,5\%) e de médico em 20 (5,2\%). A explicação é legislativa: a Portaria 3.005/2024 elevou o requisito de enfermeiro para EMAD~I de 40h (\citelink{portaria825}{Portaria 825/2016}) para 60h. Muitas equipes possuem exatamente 40h de enfermagem, configuração que era conforme pela legislação anterior. Três em cada quatro EMAD~I não conformes no país falham especificamente neste critério.

A análise regional, ilustrada na Figura~\ref{fig:conformidade_regiao}, complementa e reforça o padrão identificado na composição das equipes.

\begin{figure}[H]
\centering
\includegraphics[width=0.85\textwidth]{conformidade_por_regiao.png}
\caption{Conformidade legal por região do Brasil.}
\label{fig:conformidade_regiao}
\end{figure}

Emerge um padrão claro: as regiões com maior proporção de EMAD~I tendem a apresentar menor conformidade geral, pois a taxa de conformidade desse tipo é baixa em todas as regiões. Onde as EMAD~I representam mais da metade das equipes, puxam a média regional para baixo.

Para a modelagem do HHC-RSP, tanto a composição quanto a conformidade têm implicações diretas. A composição define uma frota heterogênea na qual cada tipo de equipe possui $Q_k$ e $S_k$ distintos; como nem toda equipe pode atender qualquer paciente, a verificação de compatibilidade $R_i \subseteq S_k$ restringe as alocações possíveis, tornando o problema mais complexo do que seria com frota homogênea. A não-conformidade agrava esse cenário: os valores reais de $Q_k$ e $S_k$ de muitas equipes ficam abaixo dos requisitos legais, reduzindo ainda mais os atendimentos viáveis. Ambos reforçam a necessidade de utilizar os dados reais de cada equipe, em vez de médias agregadas, na formulação das instâncias.

Em síntese, cobertura, composição e conformidade formam três dimensões inter-relacionadas. A distribuição desigual das equipes significa que o HHC-RSP assume escalas e restrições diferentes conforme a região: onde há poucas equipes para grandes áreas, a otimização de rotas tem maior impacto potencial. A predominância de tipos menos exigentes nas regiões de maior cobertura eleva as taxas de conformidade locais, mas a concentração de EMAD~I em outras regiões reduz a capacidade efetiva disponível para alocação. Essas diferenças, porém, se aprofundam na escala estadual, motivando uma análise mais granular.

\subsubsection{Análise na Escala Estadual}

A escala estadual é, em última instância, aquela na qual o problema de roteamento é formulado e resolvido. Ela permite identificar diferenças que a agregação regional oculta e verificar se os padrões observados se reproduzem nessa granularidade. A Figura~\ref{fig:mapa} ilustra como as equipes se distribuem espacialmente dentro de um estado, usando São Paulo como exemplo. Esse mapeamento, construído a partir das coordenadas geográficas dos estabelecimentos registradas no \citelink{cnes}{CNES}, permite identificar a localização dos depots (bases operacionais) e os recursos disponíveis para a formulação do problema de otimização em cada localidade.

\begin{figure}[H]
\centering
\includegraphics[width=0.85\textwidth]{mapaSP.png}
\caption{Distribuição espacial das equipes de Atenção Domiciliar no Estado de São Paulo, gerado com Folium (Leaflet.js). Círculos numerados são agrupamentos automáticos (clusters); as cores indicam a composição do estabelecimento (Roxo: EMAD+EMAP, Azul: apenas EMAD, Verde: apenas EMAP).}
\label{fig:mapa}
\end{figure}

\begin{figure}[H]
\centering
\includegraphics[width=0.85\textwidth]{top15_cobertura_percentual.png}
\caption{Top 15 UFs por cobertura percentual do Programa Melhor em Casa.}
\label{fig:cobertura}
\end{figure}

A Figura~\ref{fig:cobertura} evidencia a amplitude da variação entre estados: entre as quinze UFs com maior cobertura, os valores vão de 16,1\% (AM) a 100\% (DF), passando por estados como Pará (49,3\%), Rio de Janeiro (47,8\%) e Ceará (44,6\%).

\begin{table}[H]
\centering
\footnotesize
\begin{tabular}{|l|c|c|c|c|c|c|}
\hline
\textbf{Estado} & \textbf{Equipes} & \textbf{Eq./100k} & \textbf{Cob. (\%)} & \textbf{\% EMAD I} & \textbf{Conf. Geral} & \textbf{Conf. EMAD I} \\
\hline
São Paulo & 412 & 0,93 & 24,0\% & 60,9\% & 70,9\% & 59,8\% \\
Minas Gerais & 317 & 1,54 & 18,3\% & 38,8\% & 82,6\% & 82,9\% \\
Bahia & 282 & 1,99 & 33,8\% & 30,1\% & 78,7\% & 70,6\% \\
Ceará & 181 & 2,06 & 44,6\% & 40,3\% & 61,9\% & 42,5\% \\
Pará & 154 & 1,90 & 49,3\% & 37,0\% & 77,3\% & 49,1\% \\
\hline
\textbf{Brasil} & \textbf{2.664} & \textbf{1,31} & \textbf{21,9\%} & \textbf{41,0\%} & \textbf{76,7\%} & \textbf{64,6\%} \\
\hline
\end{tabular}
\caption{Indicadores dos cinco estados com mais equipes de AD.}
\end{table}

A análise estadual revela que número absoluto de equipes e cobertura relativa são indicadores dissociados: São Paulo lidera em equipes (412), mas tem a pior taxa per capita entre os cinco maiores estados, enquanto o Ceará, com menos da metade, atinge mais que o dobro por habitante. Essa dissociação se repete dentro das regiões, onde estados vizinhos como Pará e Tocantins, ou Rio de Janeiro e Espírito Santo, apresentam coberturas em extremos opostos, indicando que fatores locais como políticas de adesão e disponibilidade de profissionais são determinantes. A relação entre composição e conformidade também se mostra mais complexa na escala estadual, onde fatores locais como disponibilidade de profissionais e políticas de adesão influenciam os resultados além da composição por tipo de equipe.

Essa heterogeneidade reforça a necessidade de parametrização individualizada das instâncias do HHC-RSP, utilizando os dados reais de cada localidade em vez de médias nacionais ou regionais.

\subsubsection{Parâmetros Extraídos para o Modelo de Otimização}

A análise exploratória descrita nas seções anteriores fornece os insumos concretos para a modelagem do HHC-RSP. Foram extraídos dois parâmetros fundamentais para cada equipe: a capacidade $Q_k$ e o conjunto de habilidades $S_k$.

A capacidade $Q_k$ é definida como a soma das Cargas Horárias Semanais (CHS) de todos os profissionais da equipe:

\begin{equation}
Q_k = \sum_{p \in \text{equipe}_k} \text{CHS}_p, \quad \text{onde} \quad \text{CHS}_p = \text{CH}_{\text{amb.}} + \text{CH}_{\text{hosp.}} + \text{CH}_{\text{outros}}
\end{equation}

Nesta etapa, $Q_k$ deve ser interpretado como capacidade potencial baseada em CHS SUS; essa métrica não representa, necessariamente, dedicação exclusiva dos profissionais à Atenção Domiciliar.

O cálculo exigiu a integração de quatro tabelas do CNES: \texttt{tbEquipe} $\bowtie$ \texttt{rlEstabEquipeProf} $\bowtie$ \texttt{tbCargaHorariaSus} $\bowtie$ \texttt{tbEstabelecimento}. A distribuição resultante apresenta mediana de 230h/semana e média de 224,7h, com alta variabilidade (desvio padrão de 169,8h), refletindo a heterogeneidade entre os tipos de equipe. Os cinco estados com maior capacidade total são: São Paulo (100.929h), Minas Gerais (67.752h), Bahia (57.125h), Rio de Janeiro (51.390h) e Ceará (36.164h).

O conjunto de habilidades $S_k$ de cada equipe foi obtido mapeando códigos CBO em categorias profissionais:

\begin{equation}
S_k = \{c \in \mathcal{C} : \exists \, p \in \text{equipe}_k \text{ tal que } \text{categoria}(p) = c\}
\end{equation}

A análise identificou 17.776 profissionais únicos vinculados às equipes AD. O contingente é dominado por quatro categorias que juntas compõem 76,8\% do total: Técnico de Enfermagem (33,4\%), Enfermeiro (16,2\%), Fisioterapeuta (14,8\%) e Médico (12,4\%). As demais categorias presentes são Nutricionista, Assistente Social, Psicólogo, Fonoaudiólogo e Terapeuta Ocupacional.

Essa heterogeneidade de composição é diretamente relevante para a otimização: equipes de tipos distintos possuem capacidades ($Q_k$) e habilidades ($S_k$) diferentes, configurando um problema de frota heterogênea no qual a alocação de pacientes depende não apenas de proximidade geográfica, mas da compatibilidade $R_i \subseteq S_k$ e da capacidade residual, justificando a parametrização individualizada das instâncias.

Os parâmetros $Q_k$, $S_k$ e as coordenadas geográficas dos depots foram utilizados para gerar seis instâncias sintéticas para testes do algoritmo BRKGA, variando de 10 a 500 pacientes, com número de equipes proporcional à demanda. Como dados reais de pacientes não estão disponíveis publicamente, os pacientes foram gerados com parâmetros baseados na \citelink{portaria3005}{Portaria 3.005/2024}: modalidades AD2 (70\%) e AD3 (30\%), janelas de tempo (manhã 40\%, tarde 35\%, integral 25\%) e tempos de serviço compatíveis com a complexidade. A matriz de tempos de viagem utiliza a distância geodésica (Haversine) com velocidade média de 25~km/h.

Nesta etapa, as instâncias não incluem os requisitos de habilidades dos pacientes ($R_i$). A definição formal de $R_i$ será incorporada na próxima etapa, juntamente com as restrições de compatibilidade ($R_i \subseteq S_k$) no modelo de otimização.

\subsection{Produtos Desenvolvidos}

Todo o código-fonte está disponível publicamente no repositório GitHub: \url{https://github.com/fersuaiden/IC}. Foram desenvolvidos 11 scripts Python organizados em 4 partes (PARTE1 a PARTE4), compondo o pipeline completo de ETL e análise de dados. Os scripts geram 3 mapas interativos em HTML (mapa de equipes em SP com MarkerCluster, heatmap da capacidade no Brasil e gráfico hierárquico de composição profissional) e 39 visualizações estáticas em PNG, incluindo gráficos de cobertura, conformidade e resumo por estado. Foram também geradas 6 instâncias sintéticas para o BRKGA em formato JSON, variando de 10 a 500 pacientes, e 7 datasets processados em CSV para análises futuras. A documentação técnica (README.md) descreve o fluxo de processamento de cada etapa.

\subsection{Principais Conclusões}

A análise exploratória, ao percorrer as escalas nacional, regional e estadual, revela um cenário de cobertura insuficiente e desigual, no qual a composição das equipes e a conformidade legal variam de forma inter-relacionada. Regiões que expandiram o programa com tipos de equipe menos exigentes alcançaram maior cobertura, mas o tipo mais numeroso e operacionalmente central, a EMAD~I, apresenta baixa conformidade em todo o país, sobretudo após a elevação dos requisitos pela \citelink{portaria3005}{Portaria 3.005/2024}. Essa heterogeneidade se intensifica na escala estadual, onde estados de uma mesma região apresentam cenários operacionais muito distintos.

A caracterização completa de $Q_k$ e $S_k$, combinada com o mapeamento geográfico dos depots, fornece os insumos necessários para a modelagem do HHC-RSP. O cenário de recursos escassos, desigualmente distribuídos e parcialmente não conformes torna a otimização de rotas necessária para maximizar o impacto das equipes existentes.

\section{Estado Atual da Pesquisa}

A pesquisa encontra-se na transição entre a fase de análise de dados e a fase de modelagem. A AED das bases do CNES/DATASUS, CBO e IBGE forneceu os parâmetros de capacidade ($Q_k$), habilidades ($S_k$) e localização geográfica necessários para a formulação do HHC-RSP, além de seis instâncias sintéticas para validação inicial. A próxima etapa consiste na implementação do modelo de \citelink{mankowska2014}{Mankowska et al. (2014)} no solver Gurobi, no desenvolvimento do BRKGA com as melhorias propostas, e na comparação de desempenho entre os métodos.

\section*{Referências Bibliográficas}

\begin{thebibliography}{99}

\bibitem{ibge2022}
IBGE -- Instituto Brasileiro de Geografia e Estatística. \textit{Censo Demográfico 2022: População por idade e sexo}. Agência IBGE Notícias, 27 out. 2023. Disponível em: \url{https://agenciadenoticias.ibge.gov.br/agencia-noticias/2012-agencia-de-noticias/noticias/38186-censo-2022-numero-de-pessoas-com-65-anos-ou-mais-de-idade-cresceu-57-4-em-12-anos}. Acesso em: jan. 2026.

\bibitem{dantzig1959}
DANTZIG, G.B.; RAMSER, J.H. The truck dispatching problem. \textit{Management Science}, v. 6, n. 1, p. 80-91, 1959.

\bibitem{goncalves2011}
GONÇALVES, J.F.; RESENDE, M.G.C. Biased random-key genetic algorithms for combinatorial optimization. \textit{Journal of Heuristics}, v. 17, n. 5, p. 487-525, 2011.

\bibitem{mankowska2014}
MANKOWSKA, D.S.; MEISEL, F.; BIERWIRTH, C. The home health care routing and scheduling problem with interdependent services. \textit{Health Care Management Science}, v. 17, n. 1, p. 15-30, 2014.

\bibitem{fikar2017}
FIKAR, C.; HIRSCH, P. Home health care routing and scheduling: A review. \textit{Computers \& Operations Research}, v. 77, p. 86-95, 2017.

\bibitem{kummer2024}
KUMMER, A.; SANTOS, M.O.; ARENALES, M.N. A biased random-key genetic algorithm for the home health care routing and scheduling problem. \textit{Computers \& Operations Research}, v. 162, p. 106483, 2024.

\bibitem{portaria825}
BRASIL. Ministério da Saúde. \textit{Portaria GM/MS nº 825, de 25 de abril de 2016}. Redefine a Atenção Domiciliar no âmbito do Sistema Único de Saúde (SUS) e atualiza as equipes habilitadas. Brasília, 2016.

\bibitem{portaria3005}
BRASIL. Ministério da Saúde. \textit{Portaria GM/MS nº 3.005/2024}. Altera requisitos de composição das equipes de Atenção Domiciliar. Brasília, 2024.

\bibitem{cnes}
BRASIL. Ministério da Saúde. \textit{CNES -- Cadastro Nacional de Estabelecimentos de Saúde}. Disponível em: https://cnes.datasus.gov.br/. Acesso em: ago. 2025.

\end{thebibliography}

\end{document}
